\scanpage{14}

Une normalisation un peu plus poussée est obtenue ainsi : on se donne une trivialisation de $L$ par une section partout non nulle $D$ de $L$, [dérivation invariante], on choisit alors la base $1,x,y$ de $\mathcal E_3$ de façon que l'image de $y$ dans $L_2$ soit $i_2(D)$, celle de $x$ dans $L_3$ soit $i_3(D)$. \deleted{Cela}

\deleted{[} i.e. on prend comme structure auxiliaire la donnée d'un $\mathcal L$ \inserted{inv.}, d'un $\mathcal E$ loc. libre de rang $3$, muni d'un splitting \unsure{où} les facteurs successifs sont munis d'un isom avec $\mathcal L^0,\mathcal L^2,\mathcal L^3$ respectivement.
Le groupe d'automorphismes de cette structure est \deleted{\unsure{formé}} le sous-groupe de $G_m\times\operatorname{GL}(3)_S$, formé des \deleted{\unsure{triplets}}\inserted{couples} de la forme
\diagram{\[\left(\lambda,\left(\begin{array}{ccc}1&\alpha&\beta\\0&\lambda^2&\gamma\\0&0&\lambda^3\end{array}\right)\right)\]}
\[
\lambda\in\Gamma(G_{mS}/S)=\Gamma(S,\underline O_S^*)\qquad
\alpha,\beta,\gamma\in\Gamma(S,\underline O_S)
\]

\deleted{La fixation de la structure ainsi \unsure{conduit à une équation}} \inserted{isomorphe} \deleted{i.e.} \unsure{le} sous-groupe du \deleted{\unsure{précédent}}\inserted{groupe plein}
\diagram{\[\left(\begin{array}{ccc}1&\alpha&\beta\\0&\lambda&\gamma\\0&0&\mu\end{array}\right)\]}
pour
\[
\lambda=\left(\frac\mu\lambda\right)^2,\qquad
\mu=\left(\frac\mu\lambda\right)^3\qquad\text{i.e.}\quad\boxed{\mu^2-\lambda^3=0}.
\]

L'équation de $C$ pour une telle base \deleted{prend alors la forme} [caractérisée par la condition $\underline{x^2-y^3\in \mathcal E_5}$] prend la forme
\[
x^2-y^3+bxy+cy^2+dx+ey+f=0
\]
[i.e. on fait $a=-1$ en l'équation plus haut] \inserted{avec $\delta(b,c,d,e,f)$ une unité}
\mspage{13}
